\chapter{Graphical models and factor graphs}
\label{ch:graphical}

Chapter~\ref{ch:markov} gave a factorised density. This chapter turns factorisation into a
picture, because the picture is what makes the key result of Chapter~\ref{ch:bp} obvious
rather than surprising.

\section{Factor graphs}

A \emph{factor graph} represents a density that is a product of local terms,
\begin{equation}
\label{eq:factorised}
P(v_1,\dots,v_N) \;\propto\; \prod_{\alpha} f_\alpha\!\left(v_{\partial\alpha}\right),
\end{equation}
where each factor $f_\alpha$ depends on a subset $\partial\alpha$ of the variables. The graph
is bipartite: a circle for each variable, a square for each factor, and an edge between them
whenever the variable appears in that factor.

\begin{center}
\begin{tikzpicture}[
  var/.style={circle, draw, minimum size=7mm, inner sep=1pt},
  fac/.style={rectangle, draw, fill=black, minimum size=2.4mm, inner sep=0pt}]
\node[var] (v1) at (0,0) {$v_1$};
\node[fac] (fa) at (1.5,0) {};
\node[var] (v2) at (3,0) {$v_2$};
\node[fac] (fb) at (4.5,0) {};
\node[var] (v3) at (6,0) {$v_3$};
\draw (v1)--(fa)--(v2)--(fb)--(v3);
\node[below=1mm of fa] {\small $f_a(v_1,v_2)$};
\node[below=1mm of fb] {\small $f_b(v_2,v_3)$};
\end{tikzpicture}
\end{center}

\begin{intuition}
The graph records \emph{which variables interact directly}. Two variables not sharing a
factor influence each other only through intermediaries. That is the entire content of the
representation, and it is what will let us reorganise an intractable integral.
\end{intuition}

\section{The Markov chain as a factor graph}

For \eqref{eq:chain-factorisation}, the factors are $\mu(a_1)$ and one
$\kernel(a_i\mid a_{i-1})$ per consecutive pair. Each transition factor touches exactly two
variables, adjacent ones, so the graph is a path:

\begin{center}
\begin{tikzpicture}[
  var/.style={circle, draw, minimum size=7mm, inner sep=1pt},
  fac/.style={rectangle, draw, fill=black, minimum size=2.4mm, inner sep=0pt}]
\foreach \i in {1,2,3,4}{ \node[var] (a\i) at (2*\i,0) {$a_{\i}$}; }
\foreach \i/\j in {1/2,2/3,3/4}{
  \node[fac] (w\i) at (2*\i+1,0) {}; \draw (a\i)--(w\i)--(a\j);}
\node[fac] (m) at (1,0) {}; \draw (m)--(a1);
\node[below=1mm of m] {\small $\mu$};
\node[below=1mm of w2] {\small $\kernel$};
\end{tikzpicture}
\end{center}

\section{Trees, and why they are special}

\begin{definition}
A factor graph is a \emph{tree} if it contains no cycles: between any two nodes there is
exactly one path.
\end{definition}

A path graph is a tree. This is the structural fact everything rests on.

The reason trees are special is the \emph{separator} property.

\begin{proposition}[Separation]
\label{prop:separator}
Let the graph be a tree and let $a_i$ be a variable node. Removing $a_i$ disconnects the graph
into components. Conditional on $a_i$, the variables in different components are independent.
\end{proposition}

\begin{proof}
Because the graph is a tree, every factor involves variables lying in $\{a_i\}$ together with
exactly one component --- if a factor touched two components it would create a second path
between them, hence a cycle. Group the factors by component:
$P \propto \prod_{c} g_c\!\left(a_i, v^{(c)}\right)$ where $v^{(c)}$ are the variables of
component $c$. Fixing $a_i$, the density in the remaining variables factorises across
components, which is independence.
\end{proof}

\begin{intuition}
On a chain, this says: if you know $a_i$, then knowing anything to its left tells you nothing
further about anything to its right. All the influence between the two sides is routed through
$a_i$. That is why one can summarise everything to the left of site $i$ in a single function
of $a_i$ --- there is nothing else about the left that the right needs to know.
\end{intuition}

\begin{pitfall}
The proof used acyclicity twice. On a graph with a cycle, a factor \emph{can} join two
components, the density does not factorise on conditioning, and the summary function is no
longer sufficient. Belief propagation applied anyway (``loopy BP'') double-counts information
travelling around the cycle: a message leaving a node returns to it having been treated as
independent evidence. It may still work well in practice, but it is no longer exact, and none
of the guarantees in Chapter~\ref{ch:bp} survive.
\end{pitfall}

\section{Conditioning on the observations}
\label{sec:conditioning}

Now the step that makes the project work.

We observe $x$, related to $a$ by \eqref{eq:forward-solution}. By Bayes' rule the posterior is
prior times likelihood, and by \eqref{eq:likelihood-factorises} the likelihood factorises over
sites:
\begin{empheq}[box=\fbox]{equation}
\label{eq:posterior-factorised}
P(a\mid x,t) \;\propto\;
\underbrace{\mu(a_1)\prod_{i=2}^{\nsites}\kernel(a_i\mid a_{i-1})}_{\text{prior: chain}}
\;\cdot\;
\underbrace{\prod_{i=1}^{\nsites}\like(x_i\mid a_i)}_{\text{likelihood: unary}} .
\end{empheq}

Every likelihood factor touches \emph{one} latent variable. In factor-graph terms it is a
pendant node hanging off $a_i$ and connected to nothing else.

\begin{center}
\begin{tikzpicture}[
  var/.style={circle, draw, minimum size=7mm, inner sep=1pt},
  obs/.style={circle, draw, fill=black!12, minimum size=7mm, inner sep=1pt},
  fac/.style={rectangle, draw, fill=black, minimum size=2.4mm, inner sep=0pt}]
\foreach \i in {1,2,3,4}{
  \node[var] (a\i) at (2.2*\i,0) {$a_{\i}$};
  \node[fac] (f\i) at (2.2*\i,-1.0) {};
  \node[obs] (x\i) at (2.2*\i,-2.0) {$x_{\i}$};
  \draw (a\i)--(f\i)--(x\i);}
\foreach \i/\j in {1/2,2/3,3/4}{
  \node[fac] (w\i) at (2.2*\i+1.1,0) {}; \draw (a\i)--(w\i)--(a\j);}
\node[right=1mm of f4] {\small $\like(x_i\mid a_i)$};
\node[above=1mm of w2] {\small $\kernel(a_i\mid a_{i-1})$};
\end{tikzpicture}
\end{center}

\begin{proposition}
\label{prop:posterior-is-tree}
The factor graph of the posterior \eqref{eq:posterior-factorised} is a tree.
\end{proposition}

\begin{proof}
The latent variables and transition factors form a path, which is acyclic. Each likelihood
factor has degree one among latent variables, so attaching it adds a leaf and cannot create a
cycle. The observations are fixed, not variables.
\end{proof}

\begin{intuition}
This is the whole trick, and it is worth saying plainly why it is not obvious. Conditioning
usually \emph{creates} dependence --- explaining away is the standard example. Here it does
not, because each observation is generated from a single latent variable independently of all
others. So conditioning reweights the node potentials and leaves the edge set alone. The
posterior is a different chain, not a non-chain.
\end{intuition}

\begin{pitfall}
This fails the moment the noise is correlated across coordinates. Then
$P(x\mid a)$ does not factorise, the likelihood becomes a factor touching several $a_i$
simultaneously, and the posterior graph acquires edges --- generally becoming dense and
losing tractability entirely. The coordinatewise structure of the OU noise
(\S\ref{sec:coordinatewise}) is doing essential work here.
\end{pitfall}

\section{The cost that has been removed}

Marginalising \eqref{eq:posterior-factorised} naively to get $P(a_i\mid x,t)$ means
integrating over $\nsites-1$ variables --- exponential in the sequence length if done by brute
force on a discretisation. Proposition~\ref{prop:posterior-is-tree} says the posterior has the
structure that makes this cost linear instead. Chapter~\ref{ch:bp} carries that out.

\section*{Sources}
\addcontentsline{toc}{section}{Sources}
Factor graphs and the sum-product algorithm in the form used here are \citet{kschischang2001factor};
the graphical-model formulation and the separation properties are \citet{pearl1988}. For the
statistical-physics reading of the same objects, including exactness on trees, see
\citet{mezard2009information}; for the exponential-family and variational view, \citet{wainwright2008}.
